


\section{Experiment Details and Additional Results}\label{sec:expdetail}


\subsection{Experimental Setup}\label{sec:exp_setup}


As outlined in Section \ref{sec:exp}, we inherit most of the experimental settings from the SOTA method DyLAN \cite{dylan}, as these configurations are standard and widely validated in existing literature for direct and fair comparison.
Specifically, for all datasets, we partition the data into training, validation, and testing sets using a 3:1:1 ratio. We directly inherited most hyperparameters from DyLAN to maintain fair comparisons and, given limited computational resources, did not perform additional tuning on a validation set. The maximum token length is set to 2048 for code generation and arithmetic reasoning tasks, and 1024 for general reasoning tasks. The ranking and selection procedures for controllers optimized under Str-1 and Str-2 follow a listwise approach. To avoid positional bias, agent messages from the preceding collaboration step are randomly shuffled before being passed to agents in the subsequent step. Additionally, we employ an early-stopping mechanism to reduce unnecessary computational costs when agent outputs remain consistent across consecutive steps. Further details are available in the original DyLAN paper \cite{dylan}. For all baseline methods, we follow the implementation details and optimization configurations reported in the original papers.
All experiments are repeated three times, and the mean and standard deviation are reported.

\noindent \textbf{Experiments on general reasoning tasks.} We randomly sample 168 multiple-choice questions from the MMLU dataset \cite{mmlu} and evenly split them into training and testing sets.
The default MAS comprises seven agents with the following functional roles: ``Economist'', ``Doctor'', ``Lawyer'', ``Mathematician'', ``Psychologist'', ``Programmer'', ``Historian''. The default instruction prompts for these agents are adopted directly from DyLAN’s implementation. Answers for the agents are extracted from the agent outputs by matching the final occurrence of ``(X'' or ``(X)'', where ``X'' denotes one of the choices A, B, C, or D. The final answer is determined by selecting the option that receives the highest number of votes from agents in the last step of collaboration. The maximum number of collaboration steps is set to four, with dynamic agent selection occurring at the third step. Performance is always measured by the average classification accuracy across all questions in the four categories.



\noindent \textbf{Experiments on code generation tasks.} We sample 180 function-level code completion tasks from the HumanEval benchmark \cite{human-eval} and evenly divide them into training and testing sets.
Following DyLAN, the default MAS comprises four code-writing agents and four code-reviewing agents. Among code reviewers, three are optimizable while the fourth remains fixed as required by the method workflow.  The code writers include ``Python Assistant”, ``Algorithm Developer”, ``Computer Scientist”, and ``Programmer”. The optimizable code reviewers are ``Syntax Checker”, ``Unit Tester”, and ``Reflector”. Solutions provided by code writers undergo a review process with a maximum of six rounds. Specifically, at time steps $t = 1, 3, 4, 6$, code writers generate solutions, while at $t = 2, 5$, code reviewers provide feedback.
Dynamic agent selection occurs at the fourth step.
The final code is randomly selected from the top five code completions across all agents that pass
most tests from code reviewers. Performance evaluation is based on the average pass rate (Pass@1) of generated code across all code completion tasks.


\noindent \textbf{Experiments on arithmetic reasoning tasks.}  We randomly sample 140 mathematical problems from the MATH dataset \cite{hendrycksmath2021}, evenly covering seven subareas: algebra, counting and probability, geometry, intermediate algebra, number theory, pre-algebra, and pre-calculus. Also, these samples are split evenly into training and testing sets. In DyLAN, the authors found that collaborating agents in different domains (e.g., algebra and geometry experts) do not make significant improvement, therefore they adopt four agents with identical prompts and few-shot examples. We follow this setting for a fair comparison. The maximum number of collaboration steps is set to four, with dynamic agent selection taking place at the third step. We also follow DyLAN in using the answer extraction method described by \cite{hendrycksmath2021}. The average accuracy across all questions serves as the performance evaluation.



\subsection{Additional Experimental Results}


\subsubsection{Sensitivity of Hyper-Parameters for Single-Dimension Optimization}\label{sec:hyper_single}

% We first evaluate the performance of OMAC by varying the size of the initialized sample set (denoted as $z$) generated by the Semantic Initializer. Table \ref{tab:set_size} summarizes the results on the MATH dataset for arithmetic reasoning task. The results demonstrate that increasing the size of the initialized sample set can generally enhance the performance improvement brought by OMAC. By letting the Semantic Initializer incorporate more exploration of the design of the agent or controller being optimized, the possibility of obtaining a better agent/controller is also increasing. 

% We first evaluate the hyper-parameter sensitivity by varying the size of the initialized sample set (denoted as $z$) generated by the Semantic Initializer. Table \ref{tab:set_size} summarizes the results on the MATH dataset for arithmetic reasoning task. The results demonstrate that increasing the size of the initialized sample set can generally enhance the performance improvement brought by OMAC. By letting the Semantic Initializer incorporate more exploration of the design of the agent or controller being optimized, the possibility of obtaining a better agent/controller is also increasing. Although a larger sample size also means significant increase in computation demands because we need to evaluate each sample by conducting the  collaboration process on the entire training data, we observed a significant performance improvement against the SOTA method DyLAN with the size as mere three.


\noindent \textbf{Size of initial collection.} We first evaluate the hyper-parameter sensitivity by varying the size of initial collection (denoted as $z$) generated by the Semantic Initializer, while keeping all other settings as default. Table \ref{tab:set_size} summarizes the results on the MATH dataset for arithmetic reasoning tasks. The results demonstrate that increasing the size of the initial collection generally leads to improved performance with OMAC. This is intuitively reasonable because allowing the Semantic Initializer to explore more diverse agent/controller designs increases the likelihood of obtaining higher-performing solutions. While a larger collection entails greater computational costs since each initialized agent/controller must be evaluated via the full multi-agent collaboration on the training data, we observe substantial improvements over the SOTA DyLAN even with a modest collection size of three.




\begin{table*}[th]\footnotesize
  \centering
  \begin{threeparttable}
    \captionsetup{width=\textwidth}
    \caption{Accuracy (\%) of OMAC on each dimension with different sizes of initial collection on arithmetic reasoning tasks.}
    \label{tab:set_size}

    %%%% first (narrow) table, centered via a minipage %%%

    %%%% your full-width second table %%%%
    \begin{minipage}{0.95\textwidth}
    \begin{tabular}{%
        p{2.2cm}<{\centering}
        p{1.4cm}<{\centering}
        p{1.4cm}<{\centering}
        p{1.4cm}<{\centering}
        p{1.4cm}<{\centering}
        p{1.4cm}<{\centering}
        p{1.4cm}<{\centering}}
      \toprule
      % \multirow{2}{*}{Dimension} \\
      Collection Size & Str-1\scriptsize & Str-2\scriptsize & Str-3\scriptsize
      & Fun-1.1\scriptsize & Fun-1.2\scriptsize & Fun-2 \\
      \midrule
    $z=1$
        & 32.71\tiny{$\pm$1.63} & 32.75\tiny{$\pm$1.98}
        & 32.80\tiny{$\pm$2.32} & 33.19\tiny{$\pm$2.34}
        & 32.98\tiny{$\pm$2.82} & 32.80\tiny{$\pm$2.14} \\
    $z=2$
        & 33.14\tiny{$\pm$1.50} & 32.97\tiny{$\pm$1.41}
        & 33.25\tiny{$\pm$2.65} & 34.46\tiny{$\pm$2.40}
        & 33.90\tiny{$\pm$2.15} & 33.33\tiny{$\pm$2.49} \\
    $z=3$
        & {33.34}\tiny{$\pm$1.45} & {33.41}\tiny{$\pm$1.37}
        & {33.70}\tiny{$\pm$1.62} & {35.17}\tiny{$\pm$1.96}
        & {34.91}\tiny{$\pm$2.03} & {33.95}\tiny{$\pm$1.30} \\
    $z=4$
        & 33.69\tiny{$\pm$1.53} & 33.73\tiny{$\pm$1.42}
        & 33.96\tiny{$\pm$1.27} & 35.63\tiny{$\pm$1.52}
        & 35.21\tiny{$\pm$1.72} & 34.17\tiny{$\pm$1.53} \\
    $z=5$
        & 33.92\tiny{$\pm$1.23} & 33.96\tiny{$\pm$0.96}
        & 34.25\tiny{$\pm$1.27} & 35.89\tiny{$\pm$2.14}
        & 35.67\tiny{$\pm$1.42} & 34.51\tiny{$\pm$1.31} \\
      \bottomrule
    \end{tabular}
    \end{minipage}

  \end{threeparttable}
\end{table*}


\noindent \textbf{Maximum number of contrasting iterations.} We then conduct experiments to assess the impact of the maximum number of contrastive reasoning iterations (denoted as $w$) used by Contrastive Comparator. Table \ref{tab:iter_num} presents the results on the MATH dataset. The results indicate a similar trend: increasing the number of iterations generally results in enhanced performance by consistently refining the agents/controllers through contrastive reasoning. Furthermore, setting the iteration count to three is sufficient to achieve significant performance improvements compared to DyLAN.


\begin{table*}[ht]\footnotesize
  \centering
  \begin{threeparttable}
    \captionsetup{width=\textwidth}
    \caption{Accuracy (\%) of OMAC on each dimension with different maximum number of contrastive reasoning iterations on arithmetic reasoning tasks.}
    \label{tab:iter_num}

    %%%% first (narrow) table, centered via a minipage %%%

    %%%% your full-width second table %%%%
    \begin{minipage}{0.4795\textwidth}
    \begin{tabular}{%
        p{2.6cm}<{\centering}
        p{1.4cm}<{\centering}
        p{1.4cm}<{\centering}
        p{1.4cm}<{\centering}
        p{1.4cm}<{\centering}
        p{1.4cm}<{\centering}
        p{1.4cm}<{\centering}}
      \toprule
      % \multirow{2}{*}{Dimension} \\
      Number of Iterations & Str-1\scriptsize & Str-2\scriptsize & Str-3\scriptsize
      & Fun-1.1\scriptsize & Fun-1.2\scriptsize & Fun-2 \\
      \midrule
    $w=1$
        & 32.84\tiny{$\pm$1.23} & 32.72\tiny{$\pm$2.73}
        & 32.83\tiny{$\pm$2.41} & 33.45\tiny{$\pm$2.23}
        & 33.03\tiny{$\pm$2.31} & 32.82\tiny{$\pm$2.16} \\
    $w=2$
        & 33.02\tiny{$\pm$2.44} & 32.94\tiny{$\pm$2.14}
        & 33.07\tiny{$\pm$2.61} & 34.12\tiny{$\pm$2.24}
        & 34.12\tiny{$\pm$2.60} & 33.43\tiny{$\pm$2.57} \\
    $w=3$
        & {33.34}\tiny{$\pm$1.45} & {33.41}\tiny{$\pm$1.37}
        & {33.70}\tiny{$\pm$1.62} & {35.17}\tiny{$\pm$1.96}
        & {34.91}\tiny{$\pm$2.03} & {33.95}\tiny{$\pm$1.30} \\
    $w=4$
        & 33.85\tiny{$\pm$1.23} & 34.15\tiny{$\pm$1.14}
        & 34.41\tiny{$\pm$1.52} & 35.74\tiny{$\pm$1.12}
        & 35.26\tiny{$\pm$0.67} & 34.34\tiny{$\pm$0.93} \\
    $w=5$
        & 34.05\tiny{$\pm$0.84} & 34.77\tiny{$\pm$0.73}
        & 34.82\tiny{$\pm$1.14} & 36.04\tiny{$\pm$1.62}
        & 35.62\tiny{$\pm$1.26} & 34.79\tiny{$\pm$1.52} \\
      \bottomrule
    \end{tabular}
    \end{minipage}

    

  \end{threeparttable}
\end{table*}


\noindent \textbf{Sampling thresholds.} Lastly, We evaluate the impact of varying sampling thresholds $l$ and $h$, as described in Section \ref{sec:indi_opt}. Table \ref{tab:iter_num} presents the results for dimensions Str-1 and Fun-1.1 on the MATH dataset for arithmetic reasoning tasks. From these results, we observe that OMAC is robust to variations in the sampling thresholds, with performance fluctuations limited to within 1\%. Additionally, we note a slight performance decrease when the gap between \ $l$ and $h$ widens. This may be attributed to imbalanced sampling of positive and negative examples, potentially leading to less accurate and fair reasoning by the Contrastive Comparator.




\begin{table*}[th]\footnotesize
  \centering
  % \begin{threeparttable}
    \captionsetup{width=0.9\textwidth}
    \caption{Accuracy (\%) of OMAC on Str-1 (left) and Fun-1.1 (right) with different combinations of sampling thresholds on arithmetic reasoning tasks.}
    \label{tab:sample_thre}

    %%%% first (narrow) table, centered via a minipage %%%

    %%%% your full-width second table %%%%
    \begin{minipage}[t]{0.45\textwidth}
    \centering
    \begin{tabular}{cccc}
      \toprule
      % \multirow{2}{*}{Dimension} \\
      {} & $l=0.3$\scriptsize & $l=0.4$\scriptsize & $l=0.5$\scriptsize \\
      \midrule
    $h=0.3$
        & 32.94 & 33.02 & 32.86  \\
    $h=0.4$
        & 32.90 & {33.18} & 33.10  \\
    $h=0.5$
        & 32.79 & \textbf{33.41} & 33.34  \\
      \bottomrule
    \end{tabular}
    \end{minipage}
    \hspace{0.3cm}
    \begin{minipage}[t]{0.45\textwidth}
    \centering
    \begin{tabular}{cccc}
      \toprule
      % \multirow{2}{*}{Dimension} \\
      {} & $l=0.3$\scriptsize & $l=0.4$\scriptsize & $l=0.5$\scriptsize \\
      \midrule
    $h=0.3$
        & 34.92 & {34.95} & 34.90  \\
    $h=0.4$
        & 35.02 & 35.01 & \textbf{35.19}  \\
    $h=0.5$
        & 34.87 & 35.10 & 35.17  \\
      \bottomrule
    \end{tabular}
    \end{minipage}
    

  %  \end{threeparttable}
\end{table*}





\subsubsection{Additional Results of Ablation Study}\label{sec:add_abl}


Tables \ref{tab:abl_2} and Table \ref{tab:abl_3} summarize the results of the ablation study described in Section \ref{sec:ablation} on code generation and general reasoning tasks. The findings are consistent: the ablation model, OMAC-C, which removes the Contrastive Comparator from the optimization pipeline, performs significantly worse than the full OMAC framework. These results highlight the substantial advantage provided by the Contrastive Comparator’s contrastive reasoning on supervised positive-negative pairs.




\begin{table*}[h]\footnotesize
  \centering
  \begin{threeparttable}
    \captionsetup{width=\textwidth}
    \caption{Pass@1 (\%) of OMAC-C and OMAC with each optimization dimension on the code generation tasks.}
    \label{tab:abl_2}

    %%%% first (narrow) table, centered via a minipage %%%%
    \begin{minipage}{0.9\textwidth}
      \centering
      \begin{tabular}{cccc}
        \toprule
        Method & Str-1\scriptsize & Str-2\scriptsize & Str-3\scriptsize \\
        \midrule
        OMAC-C
          & 86.23\tiny{$\pm$1.67}
          & 85.66\tiny{$\pm$2.34}
          & 86.31\tiny{$\pm$2.80}\\
        OMAC
          & 86.76\tiny{$\pm$1.22}
          & 86.92\tiny{$\pm$2.27}
          & 87.55\tiny{$\pm$2.46}\\
        \bottomrule
      \end{tabular}
    \end{minipage}

    \vspace{0.3cm}

    %%%% your full-width second table %%%%
    \begin{minipage}{0.9\textwidth}
    \begin{tabular}{p{1.4cm}<{\centering}p{1.1cm}<{\centering}p{1.1cm}<{\centering}p{1.1cm}<{\centering}p{1.1cm}<{\centering}p{1.1cm}<{\centering}p{1.1cm}<{\centering}p{1.1cm}<{\centering}<{\centering}p{1.1cm}<{\centering}}
    \toprule
      Method & Fun-1.1\scriptsize & Fun-1.2\scriptsize & Fun-1.3\scriptsize & Fun-1.4\scriptsize & Fun-1.5\scriptsize & Fun-1.6\scriptsize & Fun-1.7\scriptsize & Fun-2\scriptsize \\
      \midrule
     OMAC-C
        & 86.74\tiny{$\pm$2.96} & 85.92\tiny{$\pm$2.42} 
        & 86.46\tiny{$\pm$2.21} & 87.86\tiny{$\pm$1.88} 
        & 87.11\tiny{$\pm$3.56} & 87.26\tiny{$\pm$2.04} 
        & 86.97\tiny{$\pm$1.95} & 86.01\tiny{$\pm$2.86} \\
    OMAC
        & 88.39\tiny{$\pm$2.54} & 86.31\tiny{$\pm$2.21} 
        & 88.87\tiny{$\pm$1.36} & 89.25\tiny{$\pm$1.30} 
        & 88.74\tiny{$\pm$2.67} & 88.39\tiny{$\pm$1.22} 
        & 88.34\tiny{$\pm$1.42} & 86.77\tiny{$\pm$2.43} \\
      \bottomrule
    \end{tabular}
    \end{minipage}

  \end{threeparttable}
\end{table*}




\endinput





\begin{table*}[h]\footnotesize
  \centering
  \begin{threeparttable}
    \captionsetup{width=\textwidth}
    \caption{Accuracy (\%) of OMAC-C and OMAC with each optimization dimension on the general reasoning tasks.}
    \label{tab:abl_3}
    %%%% first (narrow) table, centered via a minipage %%%%
    \begin{minipage}{0.9\textwidth}
      \centering
      \begin{tabular}{cccc}
        \toprule
        Method & Str-1\scriptsize & Str-2\scriptsize & Str-3\scriptsize \\
        \midrule
        OMAC-C
          & 71.13\tiny{$\pm$1.94}
          & 69.86\tiny{$\pm$1.77}
          & 70.71\tiny{$\pm$1.75}\\
        OMAC
          & 73.14\tiny{$\pm$2.24}
          & 72.06\tiny{$\pm$1.96}
          & 73.18\tiny{$\pm$1.47}\\
        \bottomrule
      \end{tabular}
    \end{minipage}

    \vspace{0.3cm}

    %%%% your full-width second table %%%%
    \begin{minipage}{0.9\textwidth}
    \begin{tabular}{p{1.4cm}<{\centering}p{1.1cm}<{\centering}p{1.1cm}<{\centering}p{1.1cm}<{\centering}p{1.1cm}<{\centering}p{1.1cm}<{\centering}p{1.1cm}<{\centering}p{1.1cm}<{\centering}<{\centering}p{1.1cm}<{\centering}}
      \toprule
      Method & Fun-1.1\scriptsize & Fun-1.2\scriptsize & Fun-1.3\scriptsize & Fun-1.4\scriptsize & Fun-1.5\scriptsize & Fun-1.6\scriptsize & Fun-1.7\scriptsize & Fun-2\scriptsize \\
      \midrule
     OMAC-C
        & 70.88\tiny{$\pm$2.55} & 70.47\tiny{$\pm$2.04} 
        & 70.46\tiny{$\pm$2.74} & 71.70\tiny{$\pm$2.15} 
        & 71.33\tiny{$\pm$3.10} & 70.19\tiny{$\pm$2.35} 
        & 70.23\tiny{$\pm$2.05} & 69.74\tiny{$\pm$2.93} \\
    OMAC
        & 72.33\tiny{$\pm$2.47} & 73.23\tiny{$\pm$1.83} 
        & 72.06\tiny{$\pm$2.41} & 74.22\tiny{$\pm$2.22} 
        & 73.15\tiny{$\pm$2.86} & 72.07\tiny{$\pm$2.92} 
        & 71.83\tiny{$\pm$2.74} & 71.02\tiny{$\pm$2.57} \\
      \bottomrule
    \end{tabular}
    \end{minipage}

  \end{threeparttable}
\end{table*}




\endinput






\subsubsection{Additional Results of Multi-Dimension Optimization}\label{sec:add_multi}

\noindent \textbf{Results on other tasks.} We present experimental results of multi-dimension optimization using our OMAC on code generation tasks and general reasoning tasks in Figure \ref{fig:mul_dim_2} and Figure \ref{fig:mul_dim_3} respectively. The trends are similar to the results shown in Section \ref{sec:multi_results}, which demonstrate that iterative optimizing multiple dimensions can bring in \textbf{significant performance improvements compared to optimizing a single dimension}. Also, only optimizing dimensions that individually demonstrate the most significant improvements is an effective strategy to bring in significant performance improvement considering the computation resource constraints.



\begin{figure}[ht]
\centering
\begin{minipage}{0.46\linewidth}
    \centering
    \includegraphics[width=0.9\linewidth]{figs/multi_3.png}
\end{minipage}
\begin{minipage}{0.46\linewidth}
    \centering
    \includegraphics[width=0.9\linewidth]{figs/multi_4.png}
\end{minipage}
%\qquad
%让图片换行，
\caption{Performance during iterative optimization for multiple dimensions on code generation tasks. The X-axis represents the dimensions undergoing three iterations of optimization, while each point indicates the MAS's performance with the optimized dimensions on the test set. The error bar denotes the standard deviation.}
% Each figure includes two dimensions optimized over three iterations, resulting in six data points in total.}
\label{fig:mul_dim_2}
% \vspace{-3mm}
\end{figure}




\begin{figure}[ht]
\centering
\begin{minipage}{0.46\linewidth}
    \centering
    \includegraphics[width=0.9\linewidth]{figs/multi_5.png}
\end{minipage}
\begin{minipage}{0.46\linewidth}
    \centering
    \includegraphics[width=0.9\linewidth]{figs/multi_6.png}
\end{minipage}
%\qquad
%让图片换行，
\caption{Performance during iterative optimization for multiple dimensions on general reasoning tasks. The X-axis represents the dimensions undergoing three iterations of optimization, while each point indicates the MAS's performance with the optimized dimensions on the test set. The error bar denotes the standard deviation.}
% Each figure includes two dimensions optimized over three iterations, resulting in six data points in total.}
\label{fig:mul_dim_3}
% \vspace{-3mm}
\end{figure}


\noindent \textbf{Validation of iterative optimization.} We further conduct experiments to validate our iterative optimization design when jointly optimizing multiple dimensions. Specifically, we propose to optimize one dimension at a time while keeping other dimensions fixed, thereby preserving the effectiveness of contrastive reasoning by limiting variability within positive-negative pairs. To verify this, we conduct experiments where two dimensions are simultaneously varied during optimization, meaning the 
Contrastive Comparator generates two agents/controllers based on the positive-negative pair in each iteration.




\begin{figure}[t]
\centering
\begin{minipage}{0.46\linewidth}
    \centering
    \includegraphics[width=0.9\linewidth]{figs/multi_7.png}
\end{minipage}
\begin{minipage}{0.46\linewidth}
    \centering
    \includegraphics[width=0.9\linewidth]{figs/multi_8.png}
\end{minipage}
%\qquad
%让图片换行，
\caption{Performance of simultaneously optimizing multiple dimensions on arithmetic reasoning tasks. The X-axis denotes the iteration number of contrastive reasoning, and each point indicates the MAS performance achieved using the optimized dimensions generated by the Contrastive Comparator. The left figure illustrates results from jointly optimizing Fun-1.1 and Fun-1.2, while the right figure corresponds to jointly optimizing Fun-1.1 and Str-3. The error bar denotes the standard deviation.}
% Each figure includes two dimensions optimized over three iterations, resulting in six data points in total.}
\label{fig:mul_dim_4}
% \vspace{-3mm}
\end{figure}

By comparing Figure \ref{fig:mul_dim_4} and Figure \ref{fig:mul_dim}, it is evident that simultaneously optimizing multiple dimensions, which asks the Contrastive Comparator to reason over multiple variable factors at once, results in significantly reduced performance gains and larger variance for OMAC. These findings validate the rationale behind our iterative multi-dimensional optimization design.




\subsubsection{Computation Cost}\label{sec:computation_apd}


% First, OMAC benefits the important online inference computation that matters for practical usage of agent systems through its efficient and effective structure control. Specifically, OMAC helps optimize to dynamically select only the most important and valuable agents for participation at certain turn of collaboration (Str-1 and Str-2), as well as intelligently determining the input context by only incorporating the valuble information. Both optimization significantly decreases the requirements of token consumption and API calls. We report the average API calls and token consumption of resolving a single problem on the code generation task (HumanEval benchmark) and arithmetic reasoning task (MATH dataset) in Table~\ref{}. From the results, we can see that OMAC can significantly decrease the computation cost for inference by dynamically optimizing the agent collaboration structure compared to other multi-agent frameworks.



First, OMAC improves online inference efficiency, which is a key factor for the practical deployment of agent systems, through its effective structural control. Specifically, OMAC optimizes to dynamically select only the most important and valuable agents to participate in the collaboration (Str-1 and Str-2) and intelligently determines the input context by incorporating only relevant information (Str-3). These optimizations substantially reduce token consumption and API calls. We report the average number of API calls and token cost with GPT-3.5-Turbo required to solve a single problem on the code generation task (HumanEval benchmark) and the arithmetic reasoning task (MATH dataset), shown in Table~\ref{tab:comput_1} and Table~\ref{tab:comput_2}. The results demonstrate that OMAC significantly reduces inference time computational cost by dynamically optimizing the collaboration structure, outperforming other multi-agent frameworks.






\begin{table*}[ht]\footnotesize
    \centering
    %\setlength{\abovecaptionskip}{3pt}
    \begin{threeparttable}
    \caption{Inference cost on code generation tasks.}
    %: RCF outperforms all baselines on all metrics with significance level $p$-value < 0.05 (indicated by $$). The last row shows the $p$-value of comparing RCF with the best baseline on the corresponding metric (indicated by boldface).}\vspace{-5pt}
    \label{tab:comput_1}
    \begin{tabular}{p{1.4cm}p{0.9cm}<{\centering}p{0.9cm}<{\centering}p{0.9cm}<{\centering}p{0.9cm}<{\centering}p{0.9cm}<{\centering}p{0.9cm}<{\centering}p{0.9cm}<{\centering}p{1.0cm}<{\centering}}
    \toprule
    \multirow{2}{*}{Method}&\multicolumn{4}{c}{Baselines}&\multicolumn{4}{c}{OMAC}\cr
    \cmidrule(lr){2-5} \cmidrule(lr){6-9}
    &\multicolumn{1}{c}{CodeT\scriptsize}&\multicolumn{1}{c}{AgentVerse\scriptsize}&\multicolumn{1}{c}{DyLAN\scriptsize}&\multicolumn{1}{c}{AFlow\scriptsize}&\multicolumn{1}{c}{Str-1\scriptsize}&\multicolumn{1}{c}{Str-2\scriptsize}&\multicolumn{1}{c}{Str-3\scriptsize}&\multicolumn{1}{c}{Str-2+Str-3\scriptsize}\cr
    \midrule
    \#API Calls & 20.72 & 23.15 & 17.94 & 17.25 & 16.34 & 15.76 & 17.03 & 14.80 \\
    Cost (\$) & 0.1430 & 0.1958 & 0.1288 & 0.1134 & 0.1032 & 0.0973 & 0.0920 & 0.0883 \\
    
    \bottomrule
    \end{tabular}

    
    \end{threeparttable}
\end{table*}




\endinput


&\multicolumn{1}{c}{SE\scriptsize}&\multicolumn{1}{c}{LLM-Blender\scriptsize}&\multicolumn{1}{c}{LLM Debate\scriptsize}&\multicolumn{1}{c}{DyLAN\scriptsize}&\multicolumn{1}{c}{Str-1\scriptsize}&\multicolumn{1}{c}{Str-2\scriptsize}&\multicolumn{1}{c}{Str-3\scriptsize}\cr






\begin{table*}[ht]\footnotesize
    \centering
    %\setlength{\abovecaptionskip}{3pt}
    \begin{threeparttable}
    \caption{Inference cost on arithmetic reasoning tasks.}
    %: RCF outperforms all baselines on all metrics with significance level $p$-value < 0.05 (indicated by $$). The last row shows the $p$-value of comparing RCF with the best baseline on the corresponding metric (indicated by boldface).}\vspace{-5pt}
    \label{tab:comput_2}
    \begin{tabular}{p{1.4cm}p{0.9cm}<{\centering}p{0.9cm}<{\centering}p{0.9cm}<{\centering}p{0.9cm}<{\centering}p{0.9cm}<{\centering}p{0.9cm}<{\centering}p{0.9cm}<{\centering}p{1.0cm}<{\centering}}
    \toprule
    \multirow{2}{*}{Method}&\multicolumn{4}{c}{Baselines}&\multicolumn{4}{c}{OMAC}\cr
    \cmidrule(lr){2-5} \cmidrule(lr){6-9}
    &\multicolumn{1}{c}{LLM-Blender\scriptsize}&\multicolumn{1}{c}{LLM-Debate\scriptsize}&\multicolumn{1}{c}{DyLAN\scriptsize}&\multicolumn{1}{c}{AFlow\scriptsize}&\multicolumn{1}{c}{Str-1\scriptsize}&\multicolumn{1}{c}{Str-2\scriptsize}&\multicolumn{1}{c}{Str-3\scriptsize}&\multicolumn{1}{c}{Str-2+Str-3\scriptsize}\cr
    \midrule
    \#API Calls & 7.12 & 8.26 & 7.30 & 7.72 & 6.85 & 6.23 & 7.12 & 5.74 \\
    Cost (\$) & 0.0458 & 0.0530 & 0.0446 & 0.0477 & 0.0369 & 0.0352 & 0.0337 & 0.0316\\
    
    \bottomrule
    \end{tabular}

    
    \end{threeparttable}
\end{table*}




\endinput


&\multicolumn{1}{c}{SE\scriptsize}&\multicolumn{1}{c}{LLM-Blender\scriptsize}&\multicolumn{1}{c}{LLM Debate\scriptsize}&\multicolumn{1}{c}{DyLAN\scriptsize}&\multicolumn{1}{c}{Str-1\scriptsize}&\multicolumn{1}{c}{Str-2\scriptsize}&\multicolumn{1}{c}{Str-3\scriptsize}\cr


For the cost on the training stage, as illustrated in Section~\ref{sec:indi_opt}, OMAC requires evaluating agents or controllers by executing the full collaborative process across the training dataset. While this approach ensures robust and representative performance evaluations,
it can increase computational costs especially when the size of training data is large and the given tasks are complex. For example, OMAC requires around 1,400 API calls for training to optimize each dimension on the HumanEval benchmark.

However, we'd like to highlight that several recent studies optimizing agent systems in a supervised and iterative manner for task-level performance, such as AVATAR~\cite{wu2024avatar} and PPDPP~\cite{deng2023plug}, also impose similar computational demands, often involving thousands of API calls during system optimization. For instance, AVATAR, which is designed to optimize only a single agent, requires over 3,000 API calls for just one training round on the FLICKR30K ENTITIES dataset \cite{plummer2015flickr30k}. Therefore, we regard these computational constraints as a common challenge across current agent system optimization methods, rather than a limitation unique to OMAC.



Nevertheless, balancing OMAC's effectiveness with computational efficiency remains a valuable and promising direction for the practical implementation of OMAC, particularly for academic researchers with limited computational resources. To investigate this, we propose a simple yet effective strategy: evaluating candidate agents or controllers on randomly sampled subsets of the training data. Although this may introduce variance, it provides a good starting point for exploring the trade-off between efficiency and robustness.

Specifically, we conducted experiments on the HumanEval benchmark, where only 50\% of the training data was randomly sampled for each evaluation. This approach reduced the number of API calls per optimization by approximately 50\%, lowering the total cost of a complete training run with GPT-3.5-turbo to around 5 dollars. Despite the reduced evaluation set, OMAC continued to outperform the strongest baseline, DyLAN, with only a minor decrease in performance compared to evaluations conducted on the full training set. These findings suggest that subset-based evaluation offers a practical and effective way to balance computational efficiency and evaluation robustness for optimization in OMAC.




\begin{table*}[ht]
  \centering
  \scriptsize
  \begin{threeparttable}
    \caption{Training cost on code generation tasks with sampled training data.}
    \label{tab:comput_3}
    % tighten spacing a bit if needed
    \setlength{\tabcolsep}{3pt}
    \begin{tabularx}{\textwidth}{
      l
      >{\centering\arraybackslash}X
      *{11}{>{\centering\arraybackslash}X}
    }
      \toprule
      \multirow{2}{*}{Method} & \multicolumn{1}{c}{Baseline} & \multicolumn{11}{c}{OMAC} \\
      \cmidrule(lr){2-2} \cmidrule(lr){3-13}
      & DyLAN
      & \mbox{Fun-1.1} & \mbox{Fun-1.2} & \mbox{Fun-1.3} & \mbox{Fun-1.4}
      & \mbox{Fun-1.5} & \mbox{Fun-1.6} & \mbox{Fun-1.7} & \mbox{Fun-2}
      & \mbox{Str-1}   & \mbox{Str-2}   & \mbox{Str-3} \\
      \midrule
      Accuracy & 85.74 & 88.27 & 86.22 & 88.62 & 89.13 & 88.51 & 88.20 & 88.23 & 86.56 & 86.43 & 86.79 & 87.47 \\
      Cost (\$)     & 4.62    & 4.95  & 5.21  & 5.02  & 5.40  & 5.67  & 5.28  & 5.15  & 4.73  & 4.62  & 5.41  & 5.24 \\
      \bottomrule
    \end{tabularx}
  \end{threeparttable}
\end{table*}






\endinput




Last, we include a comprehensive comparison on the training cost, inference cost, and performance of OMAC with the main baselines including AgentVerse, DyLAN, and AFlow. The experiments are also conducted on the HumanEval benchmark, and we report the average cost and performance across all optimization dimensions of OMAC for comparison.






\begin{table*}[ht]\footnotesize
    \centering
    %\setlength{\abovecaptionskip}{3pt}
    \begin{threeparttable}
    \captionsetup{width=0.8\textwidth}
    \caption{Efficiency comparison between OMAC's average performance and other MAS baselines on code generation tasks.}
    %: RCF outperforms all baselines on all metrics with significance level $p$-value < 0.05 (indicated by $$). The last row shows the $p$-value of comparing RCF with the best baseline on the corresponding metric (indicated by boldface).}\vspace{-5pt}
    \label{tab:efficiency}

    \begin{tabular}{ccccc}
        \toprule
        Method & AgentVerse\scriptsize & DyLAN\scriptsize & AFlow\scriptsize & OMAC (avg.)\scriptsize \\
        \midrule
        Training Cost (\$)
          & -
          & \textbf{4.62}
          & 8.24
          & 5.15\\
        Inference Cost (\$)
          & 0.1928
          & 0.1288
          & 0.1134
          & \textbf{0.0974}\\
        Accuracy (\%)
          & 78.26
          & 85.74
          & 85.63
          & \textbf{87.68}\\
        \bottomrule
        \end{tabular}


    
    \end{threeparttable}
\end{table*}




\endinput

\centering
      \begin{tabular}{cccc}
        \toprule
        Method & Str-1\scriptsize & Str-2\scriptsize & Str-3\scriptsize \\
        \midrule
        OMAC-C
          & 86.23\tiny{$\pm$1.67}
          & 85.66\tiny{$\pm$2.34}
          & 86.31\tiny{$\pm$2.80}\\
        OMAC
          & 86.76\tiny{$\pm$1.22}
          & 86.92\tiny{$\pm$2.27}
          & 87.55\tiny{$\pm$2.46}\\
        \bottomrule
      \end{tabular}

From Table~\ref{tab:efficiency}, it's clear that OMAC achieves substantial improvements in both inference efficiency and performance than all compared baselines, with only a minor increase in training cost compared to DyLAN. These findings further confirm OMAC’s strong cost–performance balance and practical applicability for real-world multi-agent systems.  



\subsubsection{Results with Different Base LLMs}\label{sec:diff_llm}

In addition to using GPT-3.5-Turbo as the base model, we also conducted experiments with three other LLMs: GPT-4o-mini and   two open-source models, DeepSeek-V2.5 and Qwen-2.5-72B. The corresponding results are reported separately in Table~\ref{tab:gpt4o}, Table~\ref{tab:deepseek}, and Table~\ref{tab:qwen}. The findings consistently demonstrate OMAC’s robustness and its sustained performance advantages over baseline methods across all optimization dimensions. These results further highlight the cross-model transferability of the optimization capabilities provided by the proposed OMAC framework.



    


\begin{table*}[ht]\footnotesize
    \centering
    %\setlength{\abovecaptionskip}{3pt}
    \begin{threeparttable}
    \caption{Performance on code generation tasks with GPT-4o-mini.}
    %: RCF outperforms all baselines on all metrics with significance level $p$-value < 0.05 (indicated by $$). The last row shows the $p$-value of comparing RCF with the best baseline on the corresponding metric (indicated by boldface).}\vspace{-5pt}
    \label{tab:gpt4o}
    \begin{tabular}{p{1.1cm}p{1.1cm}<{\centering}p{1.2cm}<{\centering}p{1.1cm}<{\centering}p{1.1cm}<{\centering}p{1.1cm}<{\centering}p{1.1cm}<{\centering}p{1.1cm}<{\centering}p{1.1cm}<{\centering}}
    \toprule
    \multirow{2}{*}{Method}&\multicolumn{5}{c}{Baselines}&\multicolumn{3}{c}{OMAC (Structural Dimension)}\cr
    \cmidrule(lr){2-6} \cmidrule(lr){7-9}
    &\multicolumn{1}{c}{CodeT\scriptsize}&\multicolumn{1}{c}{AgentVerse\scriptsize}&\multicolumn{1}{c}{ADAS\scriptsize}&\multicolumn{1}{c}{AFlow\scriptsize}&\multicolumn{1}{c}{DyLAN\scriptsize}&\multicolumn{1}{c}{Str-1\scriptsize}&\multicolumn{1}{c}{Str-2\scriptsize}&\multicolumn{1}{c}{Str-3\scriptsize}\cr
    \midrule
    Pass@1 & 79.24\tiny{$\pm$2.14} & 84.72\tiny{$\pm$1.95} & 83.83\tiny{$\pm$2.04} & 87.31\tiny{$\pm$2.70} & \underline{87.56\tiny{$\pm$2.13}} & \textbf{88.69}\tiny{$\pm$1.67} & \textbf{88.73}\tiny{$\pm$1.74} & \textbf{89.08}\tiny{$\pm$1.97} \\
    
    \bottomrule
    \end{tabular}

    \vspace{0.2cm}


    \begin{tabular}{p{1.4cm}p{1.1cm}<{\centering}p{1.1cm}<{\centering}p{1.1cm}<{\centering}p{1.1cm}<{\centering}p{1.1cm}<{\centering}p{1.1cm}<{\centering}p{1.1cm}<{\centering}<{\centering}p{1.1cm}<{\centering}}
    \toprule
    \multirow{2}{*}{Method}&\multicolumn{8}{c}{OMAC (Functional Dimension)}\cr
    \cmidrule(lr){2-9}
    &\multicolumn{1}{c}{Fun-1.1\scriptsize}&\multicolumn{1}{c}{Fun-1.2\scriptsize}&\multicolumn{1}{c}{Fun-1.3\scriptsize}&\multicolumn{1}{c}{Fun-1.4\scriptsize}&\multicolumn{1}{c}{Fun-1.5\scriptsize}&\multicolumn{1}{c}{Fun-1.6}&\multicolumn{1}{c}{Fun-1.7\scriptsize}&\multicolumn{1}{c}{Fun-2\scriptsize}\cr
    \midrule
    Pass@1 & \textbf{89.30}\tiny{$\pm$1.41} & \textbf{88.78}\tiny{$\pm$1.64} & \textbf{88.95}\tiny{$\pm$1.80} & \textbf{90.42}\tiny{$\pm$1.75} & \textbf{89.43}\tiny{$\pm$2.06} & \textbf{88.60}\tiny{$\pm$1.42} & \textbf{89.36}\tiny{$\pm$1.04} & \textbf{88.45}\tiny{$\pm$0.97}  \\
    
    \bottomrule
    \end{tabular}

    
    \end{threeparttable}
\end{table*}




\endinput






    


\begin{table*}[ht]\footnotesize
    \centering
    %\setlength{\abovecaptionskip}{3pt}
    \begin{threeparttable}
    \caption{Performance on code generation tasks with DeepSeek-V2.5.}
    %: RCF outperforms all baselines on all metrics with significance level $p$-value < 0.05 (indicated by $$). The last row shows the $p$-value of comparing RCF with the best baseline on the corresponding metric (indicated by boldface).}\vspace{-5pt}
    \label{tab:deepseek}
    \begin{tabular}{p{1.1cm}p{1.1cm}<{\centering}p{1.2cm}<{\centering}p{1.1cm}<{\centering}p{1.1cm}<{\centering}p{1.1cm}<{\centering}p{1.1cm}<{\centering}p{1.1cm}<{\centering}p{1.1cm}<{\centering}}
    \toprule
    \multirow{2}{*}{Method}&\multicolumn{5}{c}{Baselines}&\multicolumn{3}{c}{OMAC (Structural Dimension)}\cr
    \cmidrule(lr){2-6} \cmidrule(lr){7-9}
    &\multicolumn{1}{c}{CodeT\scriptsize}&\multicolumn{1}{c}{AgentVerse\scriptsize}&\multicolumn{1}{c}{ADAS\scriptsize}&\multicolumn{1}{c}{AFlow\scriptsize}&\multicolumn{1}{c}{DyLAN\scriptsize}&\multicolumn{1}{c}{Str-1\scriptsize}&\multicolumn{1}{c}{Str-2\scriptsize}&\multicolumn{1}{c}{Str-3\scriptsize}\cr
    \midrule
    Pass@1 & 79.88\tiny{$\pm$1.52} & 85.11\tiny{$\pm$2.62} & 83.98\tiny{$\pm$2.69} & \underline{87.95\tiny{$\pm$2.06}} & {87.74\tiny{$\pm$1.95}} & \textbf{89.21}\tiny{$\pm$1.76} & \textbf{88.92}\tiny{$\pm$2.31} & \textbf{89.37}\tiny{$\pm$1.47} \\
    
    \bottomrule
    \end{tabular}

    \vspace{0.2cm}


    \begin{tabular}{p{1.4cm}p{1.1cm}<{\centering}p{1.1cm}<{\centering}p{1.1cm}<{\centering}p{1.1cm}<{\centering}p{1.1cm}<{\centering}p{1.1cm}<{\centering}p{1.1cm}<{\centering}<{\centering}p{1.1cm}<{\centering}}
    \toprule
    \multirow{2}{*}{Method}&\multicolumn{8}{c}{OMAC (Functional Dimension)}\cr
    \cmidrule(lr){2-9}
    &\multicolumn{1}{c}{Fun-1.1\scriptsize}&\multicolumn{1}{c}{Fun-1.2\scriptsize}&\multicolumn{1}{c}{Fun-1.3\scriptsize}&\multicolumn{1}{c}{Fun-1.4\scriptsize}&\multicolumn{1}{c}{Fun-1.5\scriptsize}&\multicolumn{1}{c}{Fun-1.6}&\multicolumn{1}{c}{Fun-1.7\scriptsize}&\multicolumn{1}{c}{Fun-2\scriptsize}\cr
    \midrule
    Pass@1 & \textbf{89.90}\tiny{$\pm$2.52} & \textbf{88.95}\tiny{$\pm$1.92} & \textbf{89.35}\tiny{$\pm$1.52} & \textbf{90.35}\tiny{$\pm$1.47} & \textbf{89.74}\tiny{$\pm$1.67} & \textbf{89.21}\tiny{$\pm$2.04} & \textbf{89.90}\tiny{$\pm$1.81} & \textbf{89.02}\tiny{$\pm$1.74}  \\
    
    \bottomrule
    \end{tabular}

    
    \end{threeparttable}
\end{table*}




\endinput






    


\begin{table*}[h]\footnotesize
    \centering
    %\setlength{\abovecaptionskip}{3pt}
    \begin{threeparttable}
    \caption{Performance on code generation tasks with Qwen-2.5-72B.}
    %: RCF outperforms all baselines on all metrics with significance level $p$-value < 0.05 (indicated by $$). The last row shows the $p$-value of comparing RCF with the best baseline on the corresponding metric (indicated by boldface).}\vspace{-5pt}
    \label{tab:qwen}
    \begin{tabular}{p{1.1cm}p{1.1cm}<{\centering}p{1.2cm}<{\centering}p{1.1cm}<{\centering}p{1.1cm}<{\centering}p{1.1cm}<{\centering}p{1.1cm}<{\centering}p{1.1cm}<{\centering}p{1.1cm}<{\centering}}
    \toprule
    \multirow{2}{*}{Method}&\multicolumn{5}{c}{Baselines}&\multicolumn{3}{c}{OMAC (Structural Dimension)}\cr
    \cmidrule(lr){2-6} \cmidrule(lr){7-9}
    &\multicolumn{1}{c}{CodeT\scriptsize}&\multicolumn{1}{c}{AgentVerse\scriptsize}&\multicolumn{1}{c}{ADAS\scriptsize}&\multicolumn{1}{c}{AFlow\scriptsize}&\multicolumn{1}{c}{DyLAN\scriptsize}&\multicolumn{1}{c}{Str-1\scriptsize}&\multicolumn{1}{c}{Str-2\scriptsize}&\multicolumn{1}{c}{Str-3\scriptsize}\cr
    \midrule
    Pass@1 & 79.27\tiny{$\pm$2.68} & 84.14\tiny{$\pm$2.41} & 83.29\tiny{$\pm$2.67} & 86.95\tiny{$\pm$2.49} & \underline{87.02\tiny{$\pm$2.72}} & \textbf{88.32}\tiny{$\pm$1.97} & \textbf{88.56}\tiny{$\pm$2.15} & \textbf{88.91}\tiny{$\pm$1.87} \\
    
    \bottomrule
    \end{tabular}

    \vspace{0.2cm}


    \begin{tabular}{p{1.4cm}p{1.1cm}<{\centering}p{1.1cm}<{\centering}p{1.1cm}<{\centering}p{1.1cm}<{\centering}p{1.1cm}<{\centering}p{1.1cm}<{\centering}p{1.1cm}<{\centering}<{\centering}p{1.1cm}<{\centering}}
    \toprule
    \multirow{2}{*}{Method}&\multicolumn{8}{c}{OMAC (Functional Dimension)}\cr
    \cmidrule(lr){2-9}
    &\multicolumn{1}{c}{Fun-1.1\scriptsize}&\multicolumn{1}{c}{Fun-1.2\scriptsize}&\multicolumn{1}{c}{Fun-1.3\scriptsize}&\multicolumn{1}{c}{Fun-1.4\scriptsize}&\multicolumn{1}{c}{Fun-1.5\scriptsize}&\multicolumn{1}{c}{Fun-1.6}&\multicolumn{1}{c}{Fun-1.7\scriptsize}&\multicolumn{1}{c}{Fun-2\scriptsize}\cr
    \midrule
    Pass@1 & \textbf{89.26}\tiny{$\pm$1.87} & \textbf{88.83}\tiny{$\pm$2.41} & \textbf{88.74}\tiny{$\pm$1.77} & \textbf{90.02}\tiny{$\pm$1.65} & \textbf{89.13}\tiny{$\pm$1.98} & \textbf{88.64}\tiny{$\pm$1.83} & \textbf{88.95}\tiny{$\pm$2.01} & \textbf{88.16}\tiny{$\pm$1.33}  \\
    
    \bottomrule
    \end{tabular}

    
    \end{threeparttable}
\end{table*}




\endinput






\subsubsection{Experiments on Other Benchmarks}\label{sec:mbpp}

Firstly, we'd like to highlight that the benchmarks we've chosen in the main paper (HumanEval, MMLU, MATH) are well-established in the MAS literature and widely utilized by recent prominent studies such as DyLAN \cite{dylan}, ADAS \cite{hu2024automated}, LLM Debate \cite{debate1-mit}, and Agentverse \cite{chen2023agentverse}, for evaluating both multi-turn MAS and their optimization methods. Thus, we consider these benchmarks sufficiently representative to evaluate MAS performance, ensuring that our results are reliable and comparable with existing literature.

Nevertheless, to substantiate evaluate the effectiveness and generalizability of OMAC on more challenging  and complex scenarios, we conducted additional experiments on the  on the MBPP~\cite{austin2021program} benchmark and GAIA~\cite{mialon2023gaia} benchmark. MBPP benchmark consists of crowd-sourced Python programming tasks designed to emulate more complex and realistic problem-solving scenarios encountered by entry-level programmers. Similarly, we adopt DyLAN’s setup on the HumanEval benchmark as the default setting for OMAC, since both benchmarks focus on code generation tasks. As shown in Table~\ref{tab:MBPP}, the results confirm OMAC’s superior performance compared to most recent MAS optimization frameworks, including DyLAN~\cite{dylan} and AFLOW~\cite{aflow}, on more challenging tasks.


% GAIA benchmark is built with challenging real-world questions that require a set of fundamental abilities such as reasoning, multi-modality handling, web browsing, and generally tool-use proficiency to resolve. We adopt agent designs in previous research~\cite{yao2022react, huang2023agentcoder} in similar tasks to use several standard agents for tool usage as the default setting. Specifically, the default agents are Self-Consistency, Chain-of-Thought, LLM-Debate, Testing, Ensemble, Self-Refine, ReAct, respecively. The results in Table~\ref{tab:GAIA} demonstrate OMAC's superior performance over other MAS/MAS optimization approaches on the challenging tool calling benchmark.


The GAIA benchmark consists of challenging real-world questions that require a broad set of fundamental capabilities, including reasoning, multi-modal understanding, web browsing, and general tool-use proficiency. Following prior work on similar tasks~\cite{yao2022react, huang2023agentcoder}, we adopt a standard set of agents for tool-using tasks as our default configuration. Specifically, the default agent set includes Self-Consistency, Chain-of-Thought, LLM-Debate, Testing, Ensemble, Self-Refine, and ReAct. As shown in Table~\ref{tab:GAIA}, OMAC achieves superior performance compared to existing MAS and MAS optimization approaches on this challenging tool-calling benchmark.



    


\begin{table*}[ht]\footnotesize
    \centering
    %\setlength{\abovecaptionskip}{3pt}
    \begin{threeparttable}
    \caption{Performance on MBPP benchmark.}
    %: RCF outperforms all baselines on all metrics with significance level $p$-value < 0.05 (indicated by $$). The last row shows the $p$-value of comparing RCF with the best baseline on the corresponding metric (indicated by boldface).}\vspace{-5pt}
    \label{tab:MBPP}
    \begin{tabular}{p{1.4cm}p{1.3cm}<{\centering}p{1.4cm}<{\centering}p{1.4cm}<{\centering}p{1.3cm}<{\centering}p{1.2cm}<{\centering}p{1.2cm}<{\centering}p{1.2cm}<{\centering}}
    \toprule
    \multirow{2}{*}{Method}&\multicolumn{4}{c}{Baselines}&\multicolumn{3}{c}{OMAC (Structural Dimension)}\cr
    \cmidrule(lr){2-5} \cmidrule(lr){6-8}
    &\multicolumn{1}{c}{CodeT\scriptsize}&\multicolumn{1}{c}{CAMEL\scriptsize}&\multicolumn{1}{c}{Aflow\scriptsize}&\multicolumn{1}{c}{DyLAN\scriptsize}&\multicolumn{1}{c}{Str-1\scriptsize}&\multicolumn{1}{c}{Str-2\scriptsize}&\multicolumn{1}{c}{Str-3\scriptsize}\cr
    \midrule
    Pass@1(\%) & 65.72\tiny{$\pm$2.04} & 70.43\tiny{$\pm$2.41} & 81.45\tiny{$\pm$1.98} & {79.32\tiny{$\pm$2.03}} & \textbf{82.81}\tiny{$\pm$1.87} & \textbf{82.50}\tiny{$\pm$1.69} & \textbf{83.23}\tiny{$\pm$2.35} \\
    
    \bottomrule
    \end{tabular}

    \vspace{0.2cm}


    \begin{tabular}{p{1.4cm}p{1.1cm}<{\centering}p{1.1cm}<{\centering}p{1.1cm}<{\centering}p{1.1cm}<{\centering}p{1.1cm}<{\centering}p{1.1cm}<{\centering}p{1.1cm}<{\centering}<{\centering}p{1.1cm}<{\centering}}
    \toprule
    \multirow{2}{*}{Method}&\multicolumn{8}{c}{OMAC (Functional Dimension)}\cr
    \cmidrule(lr){2-9}
    &\multicolumn{1}{c}{Fun-1.1\scriptsize}&\multicolumn{1}{c}{Fun-1.2\scriptsize}&\multicolumn{1}{c}{Fun-1.3\scriptsize}&\multicolumn{1}{c}{Fun-1.4\scriptsize}&\multicolumn{1}{c}{Fun-1.5\scriptsize}&\multicolumn{1}{c}{Fun-1.6}&\multicolumn{1}{c}{Fun-1.7\scriptsize}&\multicolumn{1}{c}{Fun-2\scriptsize}\cr
    \midrule
    Pass@1(\%) & \textbf{82.49}\tiny{$\pm$2.81} & \textbf{83.30}\tiny{$\pm$2.45} &  \textbf{82.97}\tiny{$\pm$1.84} & \textbf{83.02}\tiny{$\pm$1.48} & \textbf{83.43}\tiny{$\pm$1.60} & \textbf{83.11}\tiny{$\pm$1.73} &
    \textbf{83.06}\tiny{$\pm$1.36} &
    \textbf{82.35}\tiny{$\pm$1.55} \\
    
    \bottomrule
    \end{tabular}

    
    \end{threeparttable}
\end{table*}




\endinput






    


\begin{table*}[ht]\footnotesize
    \centering
    %\setlength{\abovecaptionskip}{3pt}
    \begin{threeparttable}
    \caption{Performance on GAIA benchmark (Level 2).}
    %: RCF outperforms all baselines on all metrics with significance level $p$-value < 0.05 (indicated by $$). The last row shows the $p$-value of comparing RCF with the best baseline on the corresponding metric (indicated by boldface).}\vspace{-5pt}
    \label{tab:GAIA}
    \begin{tabular}{p{1.4cm}p{1.3cm}<{\centering}p{1.4cm}<{\centering}p{1.4cm}<{\centering}p{1.3cm}<{\centering}p{1.2cm}<{\centering}p{1.2cm}<{\centering}p{1.2cm}<{\centering}}
    \toprule
    \multirow{2}{*}{Method}&\multicolumn{4}{c}{Baselines}&\multicolumn{3}{c}{OMAC (Structural Dimension)}\cr
    \cmidrule(lr){2-5} \cmidrule(lr){6-8}
    &\multicolumn{1}{c}{AutoGPT\scriptsize}&\multicolumn{1}{c}{ADAS\scriptsize}&\multicolumn{1}{c}{Aflow\scriptsize}&\multicolumn{1}{c}{DyLAN\scriptsize}&\multicolumn{1}{c}{Str-1\scriptsize}&\multicolumn{1}{c}{Str-2\scriptsize}&\multicolumn{1}{c}{Str-3\scriptsize}\cr
    \midrule
    Pass@1(\%) & 1.04\tiny{$\pm$0.15} & 3.87\tiny{$\pm$1.25} & 8.23\tiny{$\pm$2.63} & {9.82\tiny{$\pm$2.41}} & \textbf{11.93}\tiny{$\pm$2.06} & \textbf{10.16}\tiny{$\pm$1.97} & \textbf{12.88}\tiny{$\pm$1.49} \\
    
    \bottomrule
    \end{tabular}

    \vspace{0.2cm}


    \begin{tabular}{p{1.4cm}p{1.1cm}<{\centering}p{1.1cm}<{\centering}p{1.1cm}<{\centering}p{1.1cm}<{\centering}p{1.1cm}<{\centering}p{1.1cm}<{\centering}p{1.1cm}<{\centering}<{\centering}p{1.1cm}<{\centering}}
    \toprule
    \multirow{2}{*}{Method}&\multicolumn{8}{c}{OMAC (Functional Dimension)}\cr
    \cmidrule(lr){2-9}
    &\multicolumn{1}{c}{Fun-1.1\scriptsize}&\multicolumn{1}{c}{Fun-1.2\scriptsize}&\multicolumn{1}{c}{Fun-1.3\scriptsize}&\multicolumn{1}{c}{Fun-1.4\scriptsize}&\multicolumn{1}{c}{Fun-1.5\scriptsize}&\multicolumn{1}{c}{Fun-1.6}&\multicolumn{1}{c}{Fun-1.7\scriptsize}&\multicolumn{1}{c}{Fun-2\scriptsize}\cr
    \midrule
    Pass@1(\%) & \textbf{12.14}\tiny{$\pm$2.81} & \textbf{10.92}\tiny{$\pm$2.45} &  \textbf{11.73}\tiny{$\pm$1.84} & \textbf{13.04}\tiny{$\pm$1.48} & \textbf{11.97}\tiny{$\pm$1.60} & \textbf{13.74}\tiny{$\pm$1.73} &
    \textbf{10.70}\tiny{$\pm$1.36} &
    \textbf{14.24}\tiny{$\pm$2.55} \\
    
    \bottomrule
    \end{tabular}

    
    \end{threeparttable}
\end{table*}




\endinput










\subsubsection{Training Curve}\label{sec:training_curve}


We further evaluate OMAC’s performance after each contrastive reasoning step when iteratively optimizing two dimensions (Fun1.2 then Fun1.1) over four iterations on arithmetic reasoning tasks. In particular, we report the performance curves on both the training set and the test set in Figure~\ref{fig:training_curve}.



\begin{figure}[ht]
\centering
\begin{minipage}{0.46\linewidth}
    \centering
    \includegraphics[width=0.9\linewidth]{figs/multi_9.png}
\end{minipage}
\begin{minipage}{0.46\linewidth}
    \centering
    \includegraphics[width=0.9\linewidth]{figs/multi_10.png}
\end{minipage}
%\qquad
%让图片换行，
\caption{Performance after each contrastive reasoning on arithmetic reasoning tasks. The X-axis represents the sequence number of contrastive reasoning during four iterations of optimization, while each point indicates the MAS's performance with the optimized dimensions on the training set (left) and the test set (right).}
% Each figure includes two dimensions optimized over three iterations, resulting in six data points in total.}
\label{fig:training_curve}
% \vspace{-3mm}
\end{figure}


Although OMAC’s optimization is guided by performance on a fixed training set, the results reveal a strong alignment between training and test performance, indicating that no overfitting occurs. We attribute this robustness to two key design choices. First, the optimization dimensions are constructed to target generalizable, task-level instructions rather than narrowly tailored prompts designed for specific cases. Second, each optimization iteration evaluates performance averaged across the entire training set, which inherently discourages over-specialization to particular subsets of data.

